% Vektorzeichnung: zwei Baumformen zum selben Blattwort.
\begin{center}
\setlength{\unitlength}{1pt}
\begin{picture}(270,102)
  \qbezier(50,88)(37.5,72)(25,56)
  \qbezier(50,88)(67.5,56)(85,24)
  \qbezier(25,56)(17.5,40)(10,24)
  \qbezier(25,56)(35,40)(45,24)
  \put(50,88){\circle*{3}}
  \put(25,56){\circle*{3}}
  \put(10,24){\circle*{3}}
  \put(45,24){\circle*{3}}
  \put(85,24){\circle*{3}}
  \put(10,12){\makebox(0,0){\(a\)}}
  \put(45,12){\makebox(0,0){\(b\)}}
  \put(85,12){\makebox(0,0){\(c\)}}
  \put(47,-4){\makebox(0,0){\(((ab)c)\)}}
  \qbezier(215,88)(197.5,56)(180,24)
  \qbezier(215,88)(227.5,72)(240,56)
  \qbezier(240,56)(230,40)(220,24)
  \qbezier(240,56)(247.5,40)(255,24)
  \put(215,88){\circle*{3}}
  \put(240,56){\circle*{3}}
  \put(180,24){\circle*{3}}
  \put(220,24){\circle*{3}}
  \put(255,24){\circle*{3}}
  \put(180,12){\makebox(0,0){\(a\)}}
  \put(220,12){\makebox(0,0){\(b\)}}
  \put(255,12){\makebox(0,0){\(c\)}}
  \put(218,-4){\makebox(0,0){\((a(bc))\)}}
\end{picture}
\par\medskip
\small Dasselbe Blattwort, zwei verschiedene Klammerungsbäume.
\end{center}
